\section{The intersecting lines test}\label{sec:intersection}

The introspective low-degree test forces a prover to sample a point from a register and return the evaluation of a global function at that point.
In our eventual protocol, we will want the prover to use the same global function to answer point queries sampled from \emph{multiple} different registers.
In this section, we design a game which allows us to ``transfer" global functions used from one register  to another.
We keep in mind the following picture:
\begin{equation*}
\ket{\psi} = \ket{\mathrm{EPR}^n_q}_{\reg{1}} \otimes \ket{\mathrm{EPR}^n_q}_{\reg{2}} \otimes \ket{\mathrm{aux}}_{\reg{aux}}.
\end{equation*}
We view register~$1$ as the register with the global function
and register~$2$ as the register we would like to transfer this global function to.

To accomplish this, we introduce a new test called the ``intersecting lines test".
This involves performing two introspective line-versus-point low-degree tests.
The first uses register~$1$ as its point register and register~$2$ as its slope register.
This gives us a points prover who samples~$\bu$ from register~$1$ and returns a label on it
and a line prover who samples $\bv$ from register~$2$ and returns a function on the line $\{\bu + \lambda \bv\}$,
and we know that if the points prover labels their point using a low-degree polynomial~$g$, 
then the line prover must label their line with the same polynomial~$g$.
The second low-degree test uses register~$2$ as its point register and register~$1$ as its slope register.
This gives a second line prover who returns a function on the line $\{\bv + \lambda \bu\}$.
Noting that the point $\bu + \bv$ is contained in both line provers' lines,
we can check consistency between their functions
by comparing them on this point,
forcing the second line prover to label their line using~$g$ as well.
This then entails that the second line prover from the second low-degree test must also 
label their point~$\bv$ using~$g$.
Thus, we have successfully ``transferred" the function~$g$ from the first register to the second.

In \Cref{sec:intersecting-lines}, we first introduce the intersecting lines test and prove soundness.
Following that, in \Cref{sec:introspective-lines-transfer}
we introduce an introspective version of this test which will later be used in our $\neexp$ protocol.

\subsection{The intersecting lines test}\label{sec:intersecting-lines}

\begin{definition}[Intersecting lines test]
Let $n, d > 0$ be integers, and let~$q$ be a power of~$2$.
The \emph{intersecting lines test}, denoted $\game_{\mathrm{intersect}}(n, q, d)$, is defined as follows.
Sample $\bu, \bv$ uniformly at random from $\F_q^n$,
and let $\bell$ and $\bell'$ be the two lines
$\bell := \{\bu + \lambda \bv : \lambda \in \F_q\}$
and $\bell' := \{\bv + \lambda \bu : \lambda \in \F_q\}$.
The test is performed as follows.
\begin{itemize}
\itemsep -.5pt
\item[$\circ$] The line $\bell$ and~$\bv$ are given to Alice, who responds with a degree-$d$ polynomial $\boldf: \bell \rightarrow \F_q$.
\item[$\circ$] The line $\bell'$ and~$\bu$ are given to Bob, who responds with a degree-$d$ polynomial $\boldf':\bell' \rightarrow  \F_q$.
\end{itemize}
Alice and Bob pass the test if $\boldf(\bu + \bv) = \boldf'(\bu + \bv)$.
\end{definition}

We begin by showing that although Bob knows~$\bu$, since he doesn't know~$\bv$,
the point $\bu + \bv$ looks like a uniform point in $\bell'$ to him.

\begin{fact}\label{fact:easy-peasy-lemon-peasy}
Conditioned on $\bell'$ and~$\bu$, the point~$\bu+\bv$ is distributed as a uniformly random element in $\bell'$.
\end{fact}
\begin{proof}
Let $\bw$ be a point in $\F_q^n$ such that $\bell' = \{\bw + \lambda \bu : \lambda \in \F_q\}$. 
Then for any $c \in \F_q$, $\bell'$ is also equal to the set
$\{(\bw + c\bu) + (\lambda-c) \bu : \lambda \in \F_q\}$.
Hence, $\bv$ is equally likely to be any element in the set $\{\bw + c \bu : c \in \F_q\}$,
and therefore so is $\bu + \bv$.
Since this set is also equal to $\bell'$, this proves the fact.
\end{proof}

We will be interested in the case when Alice responds using a global function~$g:\F_q^n \rightarrow \F_q$,
always setting $\boldf = g|_{\bell}$.
In this case, the following lemma shows that to succeed with high probability,
Bob must usually play the same global function as Alice.

\begin{lemma}\label{lem:intersecting-lines}
Let $(\psi, M)$ be a POVM strategy for the intersecting lines game
with value $1-\eps$.
Suppose further that there is a measurement $\{G_g\}_g$ in $\mathrm{PolyMeas}(n,d,q)$ such that
$M^{\ell, v}_{f} = G_{[g|_{\ell}=f]}$. Then
\begin{equation*}
(M^{\ell', u}_{f})_{\reg{Alice}} \otimes I_{\reg{Bob}} \consistency_{\eps + d/q} I_{\reg{Alice}} \otimes (G_{[g|_{\ell'}=f]})_{\reg{Bob}}.
\end{equation*}
\end{lemma}
\begin{proof}
Success on the test implies that
\begin{equation*}
(G_{[g(u + v) = \nu]})_{\reg{Alice}} \otimes I_{\reg{Bob}}
\consistency_{\eps} I_{\reg{Alice}} \otimes (M^{\ell', u}_{[f(u + v) = \nu]})_{\reg{Bob}}.
\end{equation*}
But by \Cref{fact:easy-peasy-lemon-peasy}, conditioned on~$\bell'$ and~$\bu$, $\bu + \bv$
is distributed as a uniformly random point in~$\bell'$.
As a result, if we let~$\bw$ be a uniformly random point in~$\bell'$, then
\begin{equation*}
(G_{[g(w) = \nu]})_{\reg{Alice}} \otimes I_{\reg{Bob}}
\consistency_{\eps} I_{\reg{Alice}} \otimes (M^{\ell', u}_{[f(w) = \nu]})_{\reg{Bob}}.
\end{equation*}
The lemma then follows from \Cref{prop:same-on-point-same-on-subspace}.
\end{proof}

\subsection{The introspective intersecting lines test}\label{sec:introspective-lines-transfer}

Now we introduce the introspective intersecting lines test.
This will be an introspective version of the intersecting lines test.

\begin{definition}
Let $n, d> 0$ be integers, let~$q$ be a power of~$2$, and let $\lambda = (2, n, q)$ be register parameters.
The \emph{introspective intersecting lines test}, denoted $\game_{\mathrm{IntroIntersect}}(\lambda, d)$,
is a $\lambda$-register game involving two registers, named ``$1$" and ``$2$", and a possible third auxiliary register.
It involves two line-versus point low-degree tests, instantiated as follows.
\begin{itemize}
\item[$\circ$] Let $\calG_1$ be a copy of $\game_{\mathrm{IntroLowDeg}}(\lambda, d)$
		using register~$1$ as the point register and register~$2$ as the slope register.
	Write $\mathsf{Lines}_1$ for the surface prover in $\calG_1$ and write $\mathsf{Points}_1$ for the points prover.
\item[$\circ$] Let $\calG_2$ be a copy of $\game_{\mathrm{IntroLowDeg}}(\lambda, d)$
		using register~$2$ as the point register and register~$1$ as the slope register.
	Write $\mathsf{Lines}_2$ for the surface prover in $\calG_2$ and write $\mathsf{Points}_2$ for the points prover.
\end{itemize}
Then $\game_{\mathrm{IntroIntersect}}(\lambda, d)$ is defined in \Cref{fig:big-transfer}.
\end{definition}

{
\floatstyle{boxed} 
\restylefloat{figure}
\begin{figure}
With probability~$\tfrac{1}{4}$ each, perform one of the following four tests.
\begin{enumerate}
	\item \textbf{Low degree test 1:} Play $\game_1$.
	\item \textbf{Low degree test 2:} Play $\game_2$.
	\item \textbf{Intersecting lines test:} Flip an unbiased coin $\bb \sim \{0, 1\}$.
		Assign the first role to Player~$\bb$ and the second role to Player~$\overline{\bb}$.
		\begin{itemize}
		\item[$\circ$] $\mathsf{Lines}_1$: Receive $\bell$, $\bv$, $\boldf:\bell\rightarrow \F_q$.
		\item[$\circ$] $\mathsf{Lines}_2$: Receive $\bell'$, $\bu$, $\boldf':\bell'\rightarrow \F_q$.
		\end{itemize}
		Accept if $\bell$ and $\bell'$ both contain $\bu + \bv$ and $\boldf(\bu + \bv) = \boldf'(\bu + \bv)$.
	\item \textbf{Consistency test:} Assign the first role to Player~$1$ and the second role to Player~$2$.
		\begin{itemize}
		\item[$\circ$] $\mathsf{Points}_1$: Receive $\bnu$.
		\item[$\circ$] $\mathsf{Points}_1$: Receive $\bnu'$.
		\end{itemize}
		Accept if $\bnu = \bnu'$.
	\end{enumerate}
	\caption{The game $\game_{\mathrm{IntroIntersect}}(\lambda, d)$.\label{fig:big-transfer}}
\end{figure}
}

\begin{remark}
We remark that although the test runs two separate introspective low-degree tests,
we cannot from these alone conclude that either of the points provers answers according to a global function.
This is because we use the lines $(k=1)$ introspective low-degree test,
whereas from \Cref{thm:big-planes-low-degree} we only know soundness for the planes $(k=2)$ introspective low-degree test.
Hence, proving soundness for the introspective intersecting lines test will require an additional assumption,
i.e.\ that one of the two points provers already answers queries according to a global function.
\end{remark}

Our main result about the introspective intersecting lines test is the following theorem.

\begin{theorem}\label{thm:big-transfer}
Let $n, d> 0$ be integers, let~$q$ be a power of~$2$, and let $\lambda = (2, n, q)$ be register parameters.
Write $\game := \game_{\mathrm{IntroIntersect}}(\lambda, d)$.
Write $A$ for the point prover's measurement in~$\game_1$,
and write~$B$ for the point prover's measurement in~$\game_2$.
\begin{itemize}
\item[$\circ$] \textbf{Completeness:}
		Suppose there is a degree-$d$ polynomial $g:\F_q^{n}\rightarrow \F_q$ such that
		\begin{equation*}
			A_{u, \nu} = \tau_u^Z \otimes I_{\reg{2}} \otimes I_{\reg{aux}} \cdot \bone[\nu = g(u)],
			\quad
			B_{v, \nu} = I_{\reg{1}} \otimes \tau_v^Z \otimes I_{\reg{aux}} \cdot \bone[\nu = g(v)].
		\end{equation*}
		Then there is a value-$1$ $\lambda$-register real commuting EPR strategy strategy for $\game$
		extending~$A$ and~$B$.
\item[$\circ$] \textbf{Soundness:}
		There exists a function $\delta(\eps) = \mathrm{poly}(\eps, d/q)$ such that the following holds.
		Let $\calS$ be a projective $\lambda$-register strategy which passes~$\game$ with probability $1-\eps$.
		Further, suppose that there exists a projective measurement
		$\{G_g\}_g$ in $\mathrm{PolyMeas}(n, d, q)$ acting on the auxiliary register such that
		\begin{equation*}
		A_{u, \nu} = \tau_u^Z \otimes I_{\reg{2}} \otimes G_{[g(u)= \nu]}.
		\end{equation*}
		Then
		\begin{equation*}
		(B_{v, \nu})_{\reg{Alice}} \otimes I_{\reg{Bob}}
			\approx_{\delta(\eps)} (I_{\reg{1}} \otimes \tau_v^Z \otimes G_{[g(v)=\nu]} )_{\reg{Alice}}\otimes I_{\reg{Bob}}.
		\end{equation*}
\end{itemize}
Furthermore,
\begin{equation*}
\qlength{\game} = O(1), \quad
\alength{\game} = O(n \log(q) + d \log(q)),
\end{equation*}
\begin{equation*}
\qtime{\game} = O(1), \quad
\atime{\game} = \poly(n \log(q), d \log(q)).
\end{equation*}
\end{theorem}


\begin{proof}[Proof of \Cref{thm:big-transfer}.]
The runtime and communication complexities follows from the $k=1$ case of the low-degree test.
The completeness follows immediately from the completeness of the introspective low-degree test.

Now, we show soundness.
Write $C$ for the line prover's measurement in~$\game_1$,
and write~$E$ for the line prover's measurement in~$\game_2$.
We can write the point provers' measurements as
\begin{equation*}
A_{u, \nu} = \tau_u^Z \otimes I_{\reg{2}} \otimes G_{[g(u)=\nu]},
\qquad B_{v, \nu} = I_{\reg{1}} \otimes \tau_v^Z \otimes B^v_\nu.
\end{equation*}
The strategy passes the consistency test with probability $1-\delta(\eps)$. As a result,
\begin{equation*}
A_{u, \nu} \otimes I_{\reg{Bob}} \consistency_{\delta(\eps)} I_{\reg{Alice}} \otimes A_{u, \nu}.
\end{equation*}
By \Cref{fact:introspective-outrospective}, this implies that
\begin{equation}\label{eq:to-use-at-the-very-end-when-you-least-expect-it}
G_{[g(u)=\nu]} \otimes I_{\reg{Bob}} \consistency_{\delta(\eps)} I_{\reg{Alice}} \otimes G_{[g(u)=\nu]}
\end{equation}
on state~$\ket{\mathrm{aux}}$ and the uniform distribution on $\F_q^n$.
By \Cref{fact:specialize-the-simeq}, this implies that
\begin{equation}\label{eq:dunno-what-to-name-this-honestly}
G_{[g|_{\ell}=f]} \otimes I_{\reg{Bob}} \consistency_{\delta(\eps)} I_{\reg{Alice}} \otimes G_{[g|_{\ell}=f]},
\end{equation}
where $\ell$ is distributed as $\bell = \{\bu + \lambda \bv : \lambda \in \F_q\}$ for uniformly random $\bu, \bv \in \F_q^n$.

Next, the strategy passes both introspective low-degree tests $\game_1$ and $\game_2$ with probability~$1-\delta(\eps)$.
By~\Cref{thm:big-planes}, this implies measurements $\{C^{\ell,v}_f\}$ and $\{E^{\ell', u}_{f'}\}$ on the auxiliary register such that
\begin{equation}\label{eq:doggone-it-one}
(C_{\ell, v, f})_{\reg{Alice}} \otimes I_{\reg{Bob}}
\approx_{\delta(\eps)} \left(\Pi^v_\ell \otimes \tau_v^Z \otimes C^{\ell,v}_f\right)_{\reg{Alice}} \otimes I_{\reg{Bob}},
\end{equation}
\begin{equation}\label{eq:doggone-it-two}
(E_{\ell', u, f'})_{\reg{Alice}} \otimes I_{\reg{Bob}}
\approx_{\delta(\eps)} \left(\tau_u^Z \otimes \Pi_{\ell'}^u \otimes E^{\ell',u}_{f'}\right)_{\reg{Alice}} \otimes I_{\reg{Bob}}.
\end{equation}
By \Cref{fact:approx-delta-game-value}, we can assume \Cref{eq:doggone-it-one} holds with equality, incurring a loss of only~$\delta(\eps)$ in the game value. (We will do the same for \Cref{eq:doggone-it-two} later.)

The strategy is now in a form that allows us to apply \Cref{lem:big-cross-check}
to the introspective cross-check in $\game_1$.
This implies that the measurements $G_{[g(u) = \nu]}$ and $C^{\ell, v}_f$ give a good strategy for the line-versus point test. In other words,
\begin{equation*}
C^{\ell,v}_{[f(u) = \nu]} \otimes I_{\reg{Bob}} \consistency_{\delta(\eps)} I_{\reg{Alice}} \otimes G_{[g(u)=\nu]}.
\end{equation*}
on state $\ket{\mathrm{aux}}$. By \Cref{prop:same-on-point-same-on-subspace}, this implies that
\begin{equation*}
C^{\ell,v}_{f} \otimes I_{\reg{Bob}} \consistency_{\delta(\eps)} I_{\reg{Alice}} \otimes G_{[g|_{\ell} = f]}.
\end{equation*}
Via~\Cref{eq:dunno-what-to-name-this-honestly}, this implies
\begin{equation*}
C^{\ell,v}_{f} \otimes I_{\reg{Bob}}
\approx_{\delta(\eps)} I_{\reg{Alice}} \otimes G_{[g|_{\ell} = f]}
\approx_{\delta(\eps)} G_{[g|_{\ell} = f]} \otimes I_{\reg{Bob}}.
\end{equation*}
As a result, by \Cref{fact:introspectivize-it},
\begin{equation*}
(C_{\ell, v, f})_{\reg{Alice}} \otimes I_{\reg{Bob}} \approx_{\delta(\eps)} (\Pi_{\ell}^v \otimes \tau_v^Z \otimes G_{[g|_{\ell}=f]})_{\reg{Alice}} \otimes I_{\reg{Bob}}.
\end{equation*}
By assumption, the right-hand side is projective.
As a result, by \Cref{fact:approx-delta-game-value}, we can assume this expression holds with equality, incurring a loss of only~$\delta(\eps)$ in the game value.
Following this, we apply \Cref{fact:approx-delta-game-value} again to assume \Cref{eq:doggone-it-two} holds with equality.

The distribution given by on $(\ell, v)$ and $(\ell', u)$ when we measure with~$C$ and~$E$ is exactly
the question distribution of the (non-introspective) intersecting lines test.
As a result, \Cref{fact:introspective-outrospective} implies that
the measurements $G_{[g_{\ell} = f]}$ and $E^{\ell', u}_{f'}$ pass the intersecting lines test with probability $1-\delta(\eps)$.
In other words,
\begin{equation*}
E^{\ell',u}_{[f'(u+v) = \nu]} \otimes I_{\reg{Alice}}
\consistency_{\delta(\eps)} I_{\reg{Bob}} \otimes G_{[g(u+v) = \nu]}.
\end{equation*}
on state $\ket{\mathrm{aux}}$.
Then by~\Cref{lem:intersecting-lines}, 
\begin{equation*}
E^{\ell',u}_{f'} \otimes I_{\reg{Bob}}
\consistency_{\delta(\eps)} I_{\reg{Alice}} \otimes G_{[g|_{\ell'} = f']}.
\end{equation*}
Via~\Cref{eq:dunno-what-to-name-this-honestly}, this implies
\begin{equation*}
E^{\ell',u}_{f'} \otimes I_{\reg{Bob}}
\approx_{\delta(\eps)} I_{\reg{Alice}} \otimes G_{[g|_{\ell} = f']}
\approx_{\delta(\eps)} G_{[g|_{\ell} = f']} \otimes I_{\reg{Bob}}.
\end{equation*}
Thus, \Cref{fact:introspectivize-it} implies that
\begin{equation*}
(E_{\ell', u, f'})_{\reg{Alice}} \otimes I_{\reg{Bob}}
\approx_{\delta(\eps)} \left(\tau_u^Z \otimes \Pi_{\ell}^u \otimes G_{[g|_{\ell'} = f']}\right)_{\reg{Alice}} \otimes I_{\reg{Bob}}.
\end{equation*}
By assumption, the right-hand side is projective.
As a result, by \Cref{fact:approx-delta-game-value}, we can assume this expression holds with equality, incurring a loss of only~$\delta(\eps)$ in the game value.

The strategy is now in a form that allows us to apply \Cref{lem:big-cross-check}
to the introspective cross-check in $\game_2$.
This implies that the measurements $B^v_{\nu}$ and $G_{[g|_{\ell'}(v) = \nu]} = G_{[g(v) = \nu]}$  give a good strategy for the line-versus-point low-degree test. In other words,
\begin{equation*}
B^v_{\nu} \otimes I_{\reg{Bob}} \consistency_{\delta(\eps)} I_{\reg{Alice}} \otimes G_{[g(v)=\nu]}.
\end{equation*}
on state $\ket{\mathrm{aux}}$.
Via~\Cref{eq:to-use-at-the-very-end-when-you-least-expect-it}, this implies
\begin{equation*}
B^{v}_{\nu} \otimes I_{\reg{Bob}}
\approx_{\delta(\eps)} I_{\reg{Alice}} \otimes G_{[g(v) = \nu]}
\approx_{\delta(\eps)} G_{[g(v) = \nu]} \otimes I_{\reg{Bob}}.
\end{equation*}
As a result, by \Cref{fact:introspectivize-it},
\begin{equation*}
(B_{v, \nu})_{\reg{Alice}} \otimes I_{\reg{Bob}} \approx_{\delta(\eps)} (I_{\reg{1}} \otimes \tau_v^Z \otimes G_{[g(v) = \nu]})_{\reg{Alice}} \otimes I_{\reg{Bob}}.
\end{equation*}
This completes the proof of the theorem.
\end{proof}
%%% Local Variables:
%%% mode: latex
%%% TeX-master: "main"
%%% End:
%%%---------- close: big-transfer

%%%%%%%%%%% jump to big-neexp-protocol
%%%---------- open: big-neexp-protocol
%!TEX root = main.tex

